\chapter{2006年安徽卷（文）}
\section{单选题}
\begin{problem}
设全集 \(U = \{1, 2, 3, 4, 5, 6, 7, 8\}\)，集合 \(S = \{1, 3, 5\}\)，\(T = \{3, 6\}\)，则 \(\complement_U(S \cup T)\) 等于（\quad）

\choices
    {\( \varnothing \)}
    {\(\{2, 4, 7, 8\}\)}
    {\(\{1,3,5,6\}\)}
    {\(\{2,4,6,8\}\)}
\end{problem}

\begin{problem}
不等式 \(\frac{1}{x} < \frac{1}{2}\) 的解集是（\quad）

\choices
    {\((-\infty, 2)\)}
    {\((2, +\infty)\)}
    {\((0, 2)\)}
    {\((-\infty, 2) \cup (2, +\infty)\)}
\end{problem}

\begin{problem}
函数 \( y = \e^{x + 1} (x \in \R) \) 的反函数是（\quad）

\choices
    {\( y = 1 + \ln x (x > 0) \)}
    {\( y = 1 - \ln x (x > 0) \)}
    {\( y = -1 - \ln x (x > 0) \)}
    {\( y = -1 + \ln x (x > 0) \)}
\end{problem}

\begin{problem}
“\(x > 3\)”是“\(x^2 > 4\)”的（\quad）

\choices
    {必要不充分条件}
    {充分不必要条件}
    {充分必要条件}
    {既不充分也不必要条件}
\end{problem}

\begin{problem}
若抛物线 \( y^{2} = 2px \) 的焦点与椭圆 \( \frac{x^2}{6} + \frac{y^2}{2} \) 的右焦点重合，则 \( p \) 的值为（\quad）

\choices
    {\(-2\)}
    {\(2\)}
    {\(-4\)}
    {\(4\)}
\end{problem}

\begin{problem}
表面积为 \(2\sqrt{3}\) 的正八面体的各个顶点都在同一球面上，则此球的体积为（\quad）

\choices
    {\(\frac{\sqrt{2}}{3}\pi\)}
    {\(\frac{1}{3}\pi\)}
    {\(\frac{2}{3}\pi\)}
    {\(\frac{2\sqrt{2}}{3}\pi\)}
\end{problem}

\begin{problem}
直线 \(x + y = 1\) 与圆 \(x^{2} + y^{2} - 2ay = 0 (a > 0)\) 没有公共点，则 \(a\) 的取值范围是（\quad）

\choices
    {\((0, \sqrt{2} - 1)\)}
    {\((\sqrt{2} - 1, \sqrt{2} + 1)\)}
    {\((- \sqrt{2} - 1, \sqrt{2} + 1)\)}
    {\((0, \sqrt{2} + 1)\)}
\end{problem}

\begin{problem}
对于函数 \( f(x) = \frac{\sin x + 1}{\sin x} (0 < x < \pi) \)，下列结论正确的是（\quad）

\choices
    {有最大值而无最小值}
    {有最小值而无最大值}
    {有最大值且有最小值}
    {既无最大值又无最小值}
\end{problem}

\begin{problem}
将函数 \( y = \sin \omega x (\omega > 0) \) 的图象按向量 \( \bs{a} = \left(-\frac{\pi}{6}, 0\right) \) 平移，平移后的图象如图所示，则平移后的图象所对应的函数解析式是（\quad）

\choices
    {\(y = \sin \left(x + \frac{\pi}{6}\right)\)}
    {\(y = \sin \left(x - \frac{\pi}{6}\right)\)}
    {\( y = \sin \left(2x + \frac{\pi}{3}\right) \)}
    {\(y = \sin \left(2x - \frac{\pi}{3}\right)\)}
\bitmapfigure[max width=0.42\linewidth,max totalheight=4.8cm]{img/2006/anhui_liberal/q9_fig1.png}
\end{problem}

\begin{problem}
如果实数 \(x, y\) 满足条件 \(\begin{cases} x - y + 1 \geqslant 0, \\ y + 1 \geqslant 0, \\ x + y + 1 \leqslant 0. \end{cases}\) 那么 \(2x - y\) 的最大值为（\quad）

\choices
    {\(2\)}
    {\(1\)}
    {\(-2\)}
    {\(-3\)}
\end{problem}

\begin{problem}
如果 \(\triangle A_{1}B_{1}C_{1}\) 的三个内角的余弦值分别等于 \(\triangle A_{2}B_{2}C_{2}\) 的三个内角的正弦值，则（\quad）

\choices
    {\(\triangle A_{1}B_{1}C_{1}\) 和 \(\triangle A_{2}B_{2}C_{2}\) 都是锐角三角形}
    {\(\triangle A_{1}B_{1}C_{1}\) 和 \(\triangle A_{2}B_{2}C_{2}\) 都是钝角三角形}
    {\(\triangle A_{1}B_{1}C_{1}\) 是钝角三角形，\(\triangle A_{2}B_{2}C_{2}\) 是锐角三角形}
    {\(\triangle A_{1}B_{1}C_{1}\) 是锐角三角形，\(\triangle A_{2}B_{2}C_{2}\) 是钝角三角形}
\end{problem}

\begin{problem}
在正方体上任选 \(3\) 个顶点连成三角形，则所得的三角形是直角非等腰三角形的概率为（\quad）

\choices
    {\(\frac{1}{7}\)}
    {\( \frac{2}{7} \)}
    {\( \frac{3}{7} \)}
    {\( \frac{4}{7} \)}
\end{problem}

\section{填空题}
\begin{problem}
设常数 \(a > 0\)，\(\left(a x^2 + \frac{1}{\sqrt{x}}\right)^4\) 展开式中 \(x^3\) 的系数为 \(\frac{3}{2}\)，则 \(a = \)\fillinblank{}.
\end{problem}

\begin{problem}
在平行四边形 \(ABCD\) 中，\(\overrightarrow{AB} = \bs{a}, \overrightarrow{AD} = \bs{b}, \overrightarrow{AN} = 3\overrightarrow{NC}, M\) 为 \(BC\) 的中点，则 \(\overrightarrow{MN} =\)\fillinblank{}（用 \(\bs{a}, \bs{b}\) 表示）.
\end{problem}

\begin{problem}
函数 \( f(x) \) 对于任意实数 \( x \) 满足条件 \( f(x + 2) = \frac{1}{f(x)} \)，若 \( f(1) = -5 \)，则 \( f(f(5)) = \)  \fillinblank{}.
\end{problem}

\begin{problem}
平行四边形的一个顶点 \(A\) 在平面 \(\alpha\) 内，其余顶点在 \(\alpha\) 的同侧，已知其中两个顶点到 \(\alpha\) 的距离分别为 \(1, 2\)，那么剩下的一个顶点到平面 \(\alpha\) 的距离可能是：

\begin{circlelist}
\item \(1\)
\item \(2\)
\item \(3\)
\item \(4\)
\end{circlelist}

以上结论正确的是 \fillinblank{}.（写出所有正确结论的编号）
\end{problem}

\section{解答题}
\begin{problem}
已知 \(0 < \alpha < \frac{\pi}{2}, \sin \alpha = \frac{4}{5}\).

\begin{enumerate}
\item 求 \( \frac{\sin^{2}\alpha + \sin 2\alpha}{\cos^{2}\alpha + \cos 2\alpha} \) 的值.
\item 求 \( \tan\left(\alpha-\frac{5\pi}{4}\right) \) 的值.
\end{enumerate}
\end{problem}

\begin{problem}
在添加剂的搭配适用中，为了找到最佳的搭配方案，需要对各种不同的搭配方式作比较，在试制某种牙膏新品种时，需要选用两种不同的添加剂. 现有芳香度分别为 \(0, 1, 2, 3, 4, 5\) 的六种添加剂可供选用. 根据试验设计学原理，通常首先要随机选取两种不同的添加剂进行搭配试验.

\begin{enumerate}
\item 求所选用的两种不同的添加剂的芳香度之和等于 \(4\) 的概率.
\item 求所选用的两种不同的添加剂的芳香度之和不小于 \(3\) 的概率.
\end{enumerate}
\end{problem}

\begin{problem}
如图，\(P\) 是边长为 \(1\) 的正六边形 \(ABCDEF\) 所在平面外一点，\(PA = 1\)，\(P\) 在平面 \(ABC\) 内的射影为 \(BF\) 的中点 \(O\).

\begin{enumerate}
\item 证明：\(PA \perp BF\).
\item 求面 \(APB\) 与面 \(DPB\) 所成二面角的大小.
\end{enumerate}

\bitmapfigure[max width=0.46\linewidth,max totalheight=5.8cm]{img/2006/anhui_liberal/q19_fig1.png}
\end{problem}

\begin{problem}
设函数 \( f(x) = x^3 + bx^2 + cx (x \in \R) \)，已知 \( g(x) = f(x) - f'(x) \) 是奇函数.

\begin{enumerate}
\item 求 \( b, c \) 的值.
\item 求 \( g(x) \) 的单调区间与极值.
\end{enumerate}
\end{problem}

\begin{problem}
在等差数列 \(\{a_{n}\}\) 中，\(a_{1} = 1\)，前 \(n\) 项和 \(S_{n}\) 满足条件 \(\frac{S_{2n}}{S_{n}} = \frac{4n + 2}{n + 1}, n = 1, 2, \cdots\).

\begin{enumerate}
\item 求数列 \(\{a_{n}\}\) 的通项公式.
\item 记 \( b_n = a_n p^{2n} (p > 0) \)，求数列 \(\{b_n\}\) 的前 \( n \) 项和 \( T_n \).
\end{enumerate}
\end{problem}

\begin{problem}
如图，\(F\) 为双曲线 \( C: \frac{x^{2}}{a^{2}} - \frac{y^{2}}{b^{2}} = 1 (a > 0, b > 0) \) 的右焦点，\(P\) 为双曲线 \(C\) 右支上一点，且位于 \(x\) 轴上方，\(M\) 为左准线上一点，\(O\) 为坐标原点. 已知四边形 \(OFPM\) 为平行四边形， \( |PF| = \lambda |OF| \).

\begin{enumerate}
\item 写出双曲线 \(C\) 的离心率 \(e\) 与 \(\lambda\) 的关系式.
\item 当 \(\lambda = 1\) 时，经过焦点 \(F\) 且平行于 \(OP\) 的直线交双曲线于 \(A\)、\(B\) 点，若 \(|AB| = 12\)，求此时的双曲线方程.
\end{enumerate}

\bitmapfigure[max width=0.46\linewidth,max totalheight=5.8cm]{img/2006/anhui_liberal/q22_fig1.png}
\end{problem}
